% main.tex -- Experience Models, ICLR 2027 submission (revised 2026-09-10 against REVIEW_HARSH_2026-09-10).
% Every number in this file is listed with its source in README.md (Number -> source table).
% No internal infrastructure identifiers, no author names, no internal repository paths.
\documentclass{article}
% TODO: the ICLR 2027 style files (iclr2027_conference.sty, iclr2027_conference.bst, math_commands.tex)
% are not present in this repository. Download them from the ICLR 2027 author kit and place them
% next to main.tex before compiling. Until they are present the file falls back to article defaults
% (the \IfFileExists branches below) so the LaTeX can be syntax-checked.
% Do NOT enable the camera-ready switch: the submission is double-blind.
\IfFileExists{iclr2027_conference.sty}{\usepackage{iclr2027_conference,times}}{\usepackage{times}\usepackage[margin=1in]{geometry}}
\usepackage[utf8]{inputenc}
\usepackage[T1]{fontenc}
\usepackage{amsmath,amssymb}
\usepackage{booktabs}
\usepackage{multirow}
\usepackage{graphicx}
\usepackage{url}
\usepackage{hyperref}
\usepackage{natbib}

\title{Experience Models: What a Per-Life Adapter\\Written From an Agent's Own Record Learns, and How It Fails}

% Anonymous for double-blind review.
\author{Anonymous authors\\Paper under double-blind review}

\newcommand{\onoff}{ON$-$OFF}
\newcommand{\nm}{\emph{not measured}}

\begin{document}
\maketitle

\begin{abstract}
We measure an \emph{experience model}: a frozen Qwen2.5-7B-Instruct agent plus a rank-8 LoRA retrained from the frozen model every 32 episodes on a success-filtered, cumulative text record of the agent's own thinking (per-life self-training; keyword retrieval over the record is present in every arm). The agent plays 1{,}024 episodes of a compiler-optimization gym and is scored every 64 episodes, adapter on and off, on 8 programs held out by identifier (the training pool contains same-source siblings of 4--6 of them); 0.49 means programs shrink 49\%. In 9 ungated lives, 7/9 life-mean gains were at or above $+0.02$ (paired SD about 0.013), but 4/9 lives produced a harmful pair (below $-0.03$), none before episode 384, every one passing a format check. A behavioural gate rejected three collapse-shaped candidates whose scores would have passed; its 0-harmful count (60 pairs) is measured on the panel it selects on (post-selection; a disjoint-panel re-probe is pending). The gain is recipe lock-in: the common adapter-on score, 0.4878, equals the birth prompt's four-pass recipe applied with no model (\texttt{-Oz}: 0.522). The adapter lowers loss on its own rows by about 3 nats per token while behaviour moves about $+0.02$. Thinking ritualizes at episodes 160--224 in 8/9 unparented lives; parenting effects ($n = 3$) are not established.
\end{abstract}

\section{Introduction}
\label{sec:intro}

An agent's experience can persist in three places: in its context, in an external text store it retrieves from, or in its weights. Most deployed language-model agents use the first two \citep{park2023generative,shinn2023reflexion,zhao2023expel}; the third is usually done once, fleet-wide, before deployment. This paper measures the third place for one agent at a time. An \emph{experience model} is a frozen instruction-tuned model plus a small per-life adapter that the agent's harness rewrites periodically from the text record of that agent's life. Mechanically it is per-life success-filtered self-training \citep{TODO-self-training}, implemented as full-rehearsal LoRA retraining from the frozen model at every ``sleep''; keyword retrieval over the same record is available in every arm, so retrieval is held constant and the quantity under test is the weight write. We keep ``experience model'' only as the artifact's name \citep{TODO-silver-sutton} and claim no novelty for its parts.

This is a characterization paper: what such a write learns on one gym over 1{,}024 episodes, how it fails, whether a behavioural gate catches the failure, and what the adapter stores versus extracts. Parenting (a stronger model prescribing thinking patterns in context) is reported as observational; a schooling corpus, a clean birth and longer lives are future work.

\paragraph{Plain words.} A \emph{chunk} is one free generation step (at most 400 tokens), scanned for the markers PREDICT, ACT, NOTE and RECALL. A \emph{probe} scores the agent on 8 programs held out from training by identifier, with fixed seeds. A \emph{paired \onoff{}} value is the probe score with the adapter loaded minus the score with it unloaded, at the same episode of the same life. A \emph{ritual} is a 32-episode window in which at least two of four repetition flags fire (same first action, flat predictions, templated notes, same recall query). The gym score is the fraction by which a program's LLVM instruction count shrinks: 0.487 means programs shrink 48.7\% on average. From the life records the adapter-OFF probe has SD 0.009 across checkpoints, so a paired value has SD about 0.013 under independence; a pair below $-0.03$ (about 2.3 paired SD) is \emph{harmful}. Where a number is not in our logged runs we write ``not measured.''

\paragraph{Contributions (measured).}
\begin{enumerate}
\item A paired, seeded probe on programs held out by identifier, with noise measured from the life records: adapter-OFF SD 0.009 (153--157 checkpoints per machine), paired SD about 0.013, same-adapter replicate SD 0.0065 with a heavy tail (one pair differing by 0.065) (\S\ref{sec:measure}).
\item Over 9 complete ungated lives, 7/9 life-mean gains at or above $+0.02$ and 0 harmful pairs through episode 320; then silent collapse in 4/9 lives (16 of 144 pairs), every harmful write passing the format canary (\S\ref{sec:safety}).
\item A score-plus-behaviour gate that rejected three collapse-shaped candidates with passing scores and five collapsing late-life candidates on score; its 0-harmful count (60 pairs, six lives, none complete) is post-selection, and it has two documented weaknesses (\S\ref{sec:safety}).
\item What the write learned is recipe lock-in: adapter-ON means repeat to four decimals across lives, and the common plateau 0.4878 equals the birth-prompt recipe applied with no model; \texttt{-Oz} scores 0.522 and \texttt{-O3} 0.341 (\S\ref{sec:storage}).
\item Storage without extraction, as a reading: about 3 nats per token of loss reduction on the adapter's own rows (0.6--1.7 above another life's rows) beside about $+0.02$ of behaviour (\S\ref{sec:storage}).
\item Ritual onset at episodes 160--224 in 8/9 measured unparented lives, recurring in 71--100\% of later windows, and present in the weights alone (\S\ref{sec:parenting}).
\item Dose and rank with their confounds: 12{,}031-row corpora pooled across lives and written once with a different trainer scored below the unseeded 0.494 base panel in 8/10 cells; frequency was not isolated; rank is unresolved at $n = 5$ (\S\ref{sec:writer}).
\item Parenting, observational ($n = 3$): onsets 192/256/288, first-NOTE overlap with the brief 0.45 against a 0.35 control null, 1 harmful pair in 48; nothing separates from controls (\S\ref{sec:parenting}).
\end{enumerate}

\paragraph{What we do not claim.} That gated writes are safe; that small frequent writes beat one large write; that rank makes no difference; that parenting improves scores, delays ritual or prevents collapse; that children take up the lesson; that the adapter stores knowledge the agent can use; anything about the schooling corpora beyond composition; anything beyond 1{,}024 episodes, eight programs, one gym and one base model. No text-memory baseline (the same agent with brief and record but no adapter) was run.

\section{The system}
\label{sec:system}

The child is Qwen2.5-7B-Instruct, frozen, plus one LoRA of rank 8 per life. A life is 1{,}024 episodes, each one program from a compiler-optimization gym played for up to 16 chunks. Constants below are design constants from the code, not measurements.

\paragraph{THINK.} The agent generates chunks of at most 400 tokens at temperature 0.7 against a context rebuilt before each call. Five line-initial markers are scanned: \texttt{PREDICT:} (expected score), \texttt{ACT:} (a pass sequence sent to the gym; the outcome is injected back), \texttt{NOTE:}, \texttt{RECALL:} (keyword retrieval over the agent's own record, 6 rows with most query-word overlap; present in every arm) and \texttt{DONE}. Every chunk, action, outcome and note is appended to a ledger written only by the harness. The one hand-written artifact is a birth prompt (Appendix~\ref{app:prompts}), identical for all arms; it names the markers, says the agent is measured on programs outside its training set, and gives one example action sequence, \texttt{-mem2reg, -sroa, -gvn, -simplifycfg}.

\paragraph{DREAM (context management).} A 22{,}000-character context: HEAD with the birth prompt, a state block and the ten most recent NOTEs; up to 6 recalled rows; a TAIL of the last 14 entries. The state block is rebuilt per episode, so NOTEs persist within an episode; across episodes only the ledger (through RECALL), a \emph{waking brief} the child writes for itself at each sleep with a verified win, and the weights carry. In parented arms the latest parent brief is also shown. This is adopted infrastructure in the style of \citet{park2023generative} and MemGPT \citep{TODO-memgpt}, not a contribution.

\paragraph{SLEEP (the write, as run).} Every 32 episodes the ledger is compiled into a corpus and a rank-8 adapter is retrained from the frozen model on the cumulative corpus. The compiler (\texttt{compile\_sleep}) is organizational: it selects \emph{wins} (actions beating the episode's running best), \emph{recoveries}, \emph{contrasts} and the 8 largest surprises; renders each in a fixed template, quoting up to two of the agent's own preceding thoughts per win; asks the child for up to six \texttt{PRINCIPLE} lines citing at least two programs (counts per sleep are logged but not tallied in our sources; \nm); strips Markdown decoration from marker lines; deduplicates on the first 160 characters. There is no paraphrase or replay-mixing step in any life. Trainer (v1, frozen): plain language-model loss on bare text, learning rate $10^{-4}$, 3 epochs, sequence length 512, $\alpha = 2r$, dropout 0.05, all seven projections, no seed flag. Nothing carries over between sleeps, so the committed adapter at any sleep is one training run over the corpus to that sleep. A trained adapter is a \emph{candidate}, committed only if it passes a format canary and, in gated arms, a probe gate (\S\ref{sec:measure}); otherwise the previous adapter stays. A second trainer (v2.1: chat template, response-only loss, learning rate $3\times10^{-5}$, 2 epochs, seed flag added 2026-09-07) was used only for offline cells (Appendix~\ref{app:writer}); a receipt audit confirmed every life used v1.

\paragraph{PARENTING (two parents; observational).} Two different parents share the name. The \emph{pattern parent} acts in the gym lives: at each sleep a ritual detector (\S\ref{sec:measure}) scores the last 32 episodes; if flagged, a stronger model (Qwen2.5-14B-Instruct, or 32B-Instruct in a 4-bit AWQ build) reads four sampled chunks and writes at most eight lines prescribing how to think, shown to the child in context. It is process-only by prompt and leak-scanned (a six-pass-name regex; since 2026-09-10 text the child itself wrote is exempt; before that, 4/16 and 6/14 briefs in two lives were replaced by a fallback). Its prompt was revised twice during the parented lives (v1 generic; v2 substance; v3 the same lesson shown at every wake, ending ``Repeat after me, every episode \ldots''; Appendix~\ref{app:prompts}); this is literal repetition, not scaffold fading. The parent's words never enter the sleep corpus; only the child's own NOTEs do. The \emph{classroom parent} teaches a separate rule game, is handed the hidden rule for diagnosis (answer-aware), and is scanned by a word-level scanner that left 5/32 delivered utterances with rule-descriptor words in one pilot round; its results are exploratory (Appendix~\ref{app:lineages}).

\paragraph{Figure 1 (placeholder).} \emph{THINK (context, chunk, markers, ledger); DREAM (HEAD/recalled/TAIL; brief across wakes); SLEEP (ledger $\to$ select $\to$ canonicalize $\to$ render $\to$ dedup $\to$ LoRA from the frozen model $\to$ canary $\to$ [gate] $\to$ commit or keep previous); PARENTING (detector $\to$ brief $\to$ leak scan $\to$ context; the brief never enters the corpus).}

\section{Measurement design}
\label{sec:measure}

\paragraph{Gym and split.} CompilerGym \texttt{llvm-v0} \citep{TODO-compilergym}, instruction-count observation; an action is an ordered list of pass names. The eight probe programs are cbench-v1 programs held out by identifier only: the training pool is the remaining 15 cbench-v1 plus 12 chstone-v0 and 40 mibench-v1 programs (67, each replayed 15--16 times per life), and it contains same-source siblings of four to six probes (mibench \texttt{dijkstra}, \texttt{gsm}, \texttt{jpeg-c}, \texttt{susan} share a name and source package with probes; training libtiff/libjpeg programs share code with \texttt{tiff2bw} and \texttt{jpeg-c}; only \texttt{sha} and \texttt{patricia} have no known sibling; Appendix~\ref{app:perprogram}). An earlier pilot leaked the probe set through an identifier prefix; the fix normalizes identifiers, and all life numbers post-date it. The 12{,}031-row rank corpus (\S\ref{sec:writer}) was pooled from lives including pre-fix ones and contains 51 rows of the probe programs, so that table is not held out.

\paragraph{Probe and paired \onoff{}.} A probe runs the eight programs from a fresh context holding only the birth prompt, up to 16 chunks each, each chunk's seed derived from the program, the chunk index and one fixed probe seed. Every 64 episodes it runs with the current adapter and with the bare model; the paired value is the difference. Probes are seeded but not common-random (sleep-time calls are unseeded; the prompt carries a wall-clock field). \emph{Noise, from the life records.} The adapter-OFF probe read 0.463--0.488 across checkpoints, mean 0.471--0.474, SD 0.0092 and 0.0087 on the two machines ($n = 157$ and 153); a paired value has SD about 0.013 under independence. Same-adapter, same-seed replicates have SD 0.0064 and 0.0067 ($n = 33$, 37) but a heavy tail: one pair differed by 0.065 (\S\ref{sec:safety}). A separate unseeded base-only script (24-chunk budget) gave $0.4936 \pm 0.0113$ over nine panels; the offline cells were counted against it. \emph{Thresholds.} Harmful $=$ pair below $-0.03$ (about 2.3 paired SD); gate tolerance 0.02 (about 1.5 paired SD, 3 replicate SD); both were set against the 0.0113 single-panel SD.

\paragraph{Cell counts, not $p$-values.} Same-configuration replicates disagreed by more than the noise band (two default-seed v2.1 replicates of one 1{,}200-row cell by 0.059, two 12{,}031-row replicates by 0.078, two unseeded panels of one adapter by 0.147); trainer nondeterminism and probe sampling were not separated. We report counts of lives, pairs and cells, treat the life as the unit, and label every inferred mechanism a reading with its falsifier.

\paragraph{Format canary and probe gate.} The canary runs the candidate on four probe programs for three chunks each and passes if at least half the chunks contain a parseable \texttt{ACT:} line. The probe gate (arms R3, R4; built after the collapses of \S\ref{sec:safety}) re-runs the candidate on the same eight programs with the same probe seed and commits only if its score is at least $\max(\text{latest OFF}, \text{latest ON}) - 0.02$ and it uses at least half the base's chunks per episode. Because the gate selects on the panel and seed that report the score, every R3/R4 paired probe is a post-selection statistic and their 0-harmful counts largely restate the gate rule. A disjoint re-probe panel (12 programs from datasets no curriculum touches) is defined, and re-probing every committed final and one mid-life adapter per life against the base is pending GPU time (Appendix~\ref{app:gates}). The gate has not been ratified by the protocol reviewer.

\paragraph{Ritual and echo.} At each sleep four flags are computed over the last 32 episodes: same recipe (modal first action in $\ge 70\%$), flat predictions (SD $< 0.05$ over $\ge 8$), templated notes (mean Jaccard of consecutive first NOTEs $\ge 0.5$), same recall ($\ge 70\%$ of $\ge 8$); ritual $=$ at least two flags. Cutoffs were fixed when the parented lives were at or below episode 48; the templated-notes rule was changed once, on 2026-09-10 for the parented arm only, to drop NOTEs overlapping the brief (Appendix~\ref{app:ritual}). The \emph{echo metric} is the fraction of episodes whose first NOTE shares at least four content words with the latest brief (stop-words and the brief's fixed closing instruction removed), scored also on unbriefed control lives against the same briefs; it is lexical, not evidence of uptake.

\paragraph{Absorption probe.} For each committed sleep $k$ we sample 24 rows from the life's cumulative corpus and 24 from another same-design life's final corpus (same gym; not a gym-free control) and take base-minus-adapter mean per-token loss on each. Rows are scored under a chat template though the v1 trainer saw bare text (rescoring queued); per-row losses were not saved, so no interval exists (2 lives). Retention is not testable under cumulative retraining, because earlier rows remain in every later training set.

\paragraph{Contamination rules.} Pass names, program names and scores are filtered from anything the parent says. Any child that has played the gym is a disposable evaluation life. The three schooling corpora are development-only (\S\ref{sec:discussion}).

\section{Results}
\label{sec:results}

Arms: R2 (fixed writer, canary only, no parent; 9 lives, complete), R3 (probe gate; 6 lives, none complete), RP (pattern parent, no gate; 3 lives, complete), R4 (gate + parent; 7 lives, early). Where an earlier dated snapshot exists, the life-level recount of 2026-09-10 is the number we stand on; snapshots are in Appendix~\ref{app:lives}.\footnote{Numbers are as of the dated recount of 2026-09-10. Proposed freeze: abstract numbers on 2026-09-14, paper numbers on 2026-09-20; any number that changes at either freeze will be updated with its README entry.}

\subsection{The writer: dose and rank, with confounds}
\label{sec:writer}

Table~\ref{tab:rankdose}: one-shot adapters written with the v2.1 trainer from corpora pooled across lives. At 300 rows the ranks tie; at 1{,}200 rows two same-configuration replicates disagree in sign (r16 $-$ r8 $= -0.036$ and $+0.009$), so an earlier ``rank 16 harms'' reading was withdrawn; at 12{,}031 rows (every unique row from all lives, including 51 probe-program rows) the adapters score below the unseeded 0.494 base panel in 8/10 cells (r8 5/5, r16 3/5; 7/10 against 0.485; against a base probed with the cells' own procedure, \nm). Paired by machine and training seed, r16 $-$ r8 was $+0.060, +0.064, -0.015, +0.011, +0.091$ (mean $+0.042$, SD 0.043, 90\% interval $[+0.001, +0.083]$); under the rule fixed before the cells ran (equivalence only if the interval lies inside $\pm 0.03$), rank is unresolved at $n = 5$; rank 8 stayed in use as the running default, not because a tie was shown. Three confounds separate these cells from the lives: trainer (v2.1 vs.\ v1), corpus provenance (pooled vs.\ a life's own successive adapters), and base reference. Because every sleep retrains from the frozen model, a life's committed adapter is itself one write over its corpus so far, so frequency was not isolated and this table does not show that small frequent writes beat one large write. Reading, held loosely: a single write of a large corpus pooled across lives scores worse than no adapter. Falsifier: 12k$\to$2k subsample cells under the same trainer and provenance land at base.

\begin{table}[t]
\centering
\small
\caption{One-shot writes from the frozen model (v2.1 trainer), probe mean on the 8 programs held out by identifier; the 12{,}031-row corpus contains 51 rows of these programs, so this table is not held out. Base $\approx 0.494$ (SD 0.0113) was measured with a different probe procedure (unseeded, 24-chunk budget) than the cells (16-chunk), so below-base counts are approximate. Each cell is one adapter scored by two unseeded panels; 1{,}200- and 12{,}031-row cells are same-configuration replicates on two machines.}
\label{tab:rankdose}
\begin{tabular}{lll}
\toprule
Rows & rank 8 & rank 16 \\
\midrule
300 & 0.485 (panels 0.477 / 0.492) & 0.488 (0.491 / 0.485) \\
1{,}200 & 0.492, 0.506 & 0.456, 0.515 \\
12{,}031 & 0.413, 0.440, 0.468, 0.491, 0.343 (mean 0.431, SD 0.057) & 0.473, 0.504, 0.453, 0.502, 0.434 (mean 0.473, SD 0.030) \\
Below 0.494 at 12{,}031 & 5/5 & 3/5 \\
\bottomrule
\end{tabular}
\end{table}

\subsection{Safety: ungated late-life collapse, then the gate}
\label{sec:safety}

\paragraph{Ungated lives (R2).} Table~\ref{tab:perlife}, top. Nine lives ran with the fixed writer and the canary as the only commit check (active: 16 rejections in seed 2, 10 in seed 0). Over complete lives (16 pairs each), 7/9 life means are at or above $+0.02$ and 2/9 negative; 4/9 lives (seeds 2, 3, 5, 6) had at least one harmful pair, 3/9 two or more, 16 harmful pairs of 144; the negative lives bottomed at $-0.242$ (seed 5) and $-0.051$ (seed 2). A dated mid-life snapshot (82 of 144 pairs: 69 positive, 9 harmful) records no harmful pair through episode 320 in any life, the first at 384, and all nine snapshot harms at or after 384 in three lives (seed 5 to $-0.186$ at 384--576; seed 6 to $-0.089$ at 448--576, about 0 by 704; seed 2 at 640--768); the recount does not record per-pair episodes, so the age of the last ungated harm over full lives is \nm. A certification draft from a still earlier snapshot (21/21 pairs positive, mean $+0.022$) had recommended freezing the writer; the later counts reversed it. Seeds 0, 1 and 5 were interrupted once by an operator error and resumed from checkpoint, re-recording partial ledger rows into later corpora; seed 5, the worst life, is among them (Appendix~\ref{app:lives}).

\paragraph{Every harmful sleep passed the format canary.} Two shapes were inspected. (i) \emph{Act once and stop}: at sleep 448 of seed 5 the adapter-ON probe used 1.0--1.8 chunks per episode with a parseable-\texttt{ACT} rate of 1.00, against 13 chunks for the base; canary 1.00 ($n = 1$). (ii) \emph{Frozen on a bad adapter}: seed 2's sleep-608 adapter passed the canary at 0.67 (threshold 0.5), the next three pairs were harmful, and every later candidate failed the canary (0.08, 0.17, 0.00, \ldots), leaving the life on the bad adapter through 16 rejections; the runner has no recovery path. Against the certification bar (absorption, retention, interface, no-harm): absorption held in the two lives probed (\S\ref{sec:storage}); retention is not testable; interface holds only as the canary measures it, an instrument shown here to miss behavioural collapse; no-harm fails from about the twelfth sleep in 4/9 lives.

\paragraph{Gated lives (R3): post-selection throughout.} Table~\ref{tab:perlife}, bottom. The gate selects on the eight programs and seed that report the score, with tolerance 0.02 tighter than the harm threshold 0.03, so ``0 harmful in 60 pairs across six lives'' is largely a selection statistic: read it as ``no committed adapter fell below its gate floor on the selection panel so far,'' not ``safe.'' The evidence that the gate refuses collapse is its rejection log. Three candidates with passing scores and perfect format were rejected for brevity before commit: life 500 at sleep 64 (0.4878 vs.\ floor 0.4847; 2.1 vs.\ 11.6 chunks per episode), life 505 at sleep 288 (0.520 vs.\ 0.519; 1.75 vs.\ 9.6), life 501 at sleep 704 (0.488 at floor; 4.4 vs.\ 11.4). In life 500, five consecutive late-life candidates (sleeps 672--800) scored 0.194--0.341 against a floor of 0.527--0.529 with canary 1.00 and 6.6--10 chunks per episode and were rejected on score; they were collapsing as ungated seed 5's writes had. Exposure: four of six lives have passed episode 384 and three have passed 768; none has completed. Two weaknesses are on record. First, the 0.02 slack also applies against the previous adapter, so a candidate up to 0.02 below the adapter it replaces still replaces it (observed once: 0.4903 replaced 0.5065). Second, the same committed adapter in life 500 read 0.5291 at five consecutive probes (episodes 448--704) and 0.4638 at 768 with no commit between; same-adapter replicates usually agree (SD 0.0065 over 70 pairs) but this pair differs by 0.065, so one 8-episode probe can misjudge an adapter by about 0.06 on a rare draw, and because the floor is the latest ON probe that draw lowered the floor from 0.529 to 0.464. A confirm-at-next-probe rule and a best-confirmed-probe floor are proposed, not built.

\begin{table}[t]
\centering
\small
\caption{Per-life table from the probe and gate records (recount of 2026-09-10; unit $=$ life). Paired \onoff{} on the 8 programs held out by identifier; harmful $=$ pair $< -0.03$ (paired SD $\approx 0.013$). Exposure $=$ episodes reached; restarts from the life logs. R3 and R4 rows are post-selection (the gate selects on the panel and seed that report the score; disjoint-panel re-probe pending) and are not comparable to R2/RP means. Earlier dated snapshots are in Appendix~\ref{app:lives}.}
\label{tab:perlife}
\begin{tabular}{llrrrrrrl}
\toprule
Life & Arm & Pairs & Mean & Min & Harmful & Exposure & Restarts & Commits / rejections \\
\midrule
R2 seed 0 & ungated & 16 & $+0.037$ & $-0.000$ & 0 & 1{,}024 & 1 & 22 / canary 10 \\
R2 seed 1 & ungated & 16 & $+0.050$ & $+0.003$ & 0 & 1{,}024 & 1 & 32 / 0 \\
R2 seed 2 & ungated & 16 & $-0.000$ & $-0.051$ & 5 & 1{,}024 & 0 & 16 / canary 16 \\
R2 seed 3 & ungated & 16 & $+0.045$ & $-0.057$ & 1 & 1{,}024 & 0 & 32 / 0 \\
R2 seed 4 & ungated & 16 & $+0.022$ & $+0.009$ & 0 & 1{,}024 & 0 & 30 / canary 2 \\
R2 seed 5 & ungated & 16 & $-0.051$ & $-0.242$ & 8 & 1{,}024 & 1 & 31 / canary 1 \\
R2 seed 6 & ungated & 16 & $-0.005$ & $-0.089$ & 2 & 1{,}024 & 0 & 32 / 0 \\
R2 seed 7 & ungated & 16 & $+0.050$ & $+0.023$ & 0 & 1{,}024 & 0 & 32 / 0 \\
R2 seed 8 & ungated & 16 & $+0.028$ & $+0.000$ & 0 & 1{,}024 & 0 & 30 / canary 2 \\
RP 400 & pattern parent (14B) & 16 & $+0.038$ & $+0.006$ & 0 & 1{,}024 & 3 & 32 / 0 \\
RP 401 & pattern parent (14B) & 16 & $+0.016$ & $-0.005$ & 0 & 1{,}024 & 3 & 32 / 0 \\
RP 402 & pattern parent (32B-AWQ) & 16 & $+0.011$ & $-0.039$ & 1 & 1{,}024 & 4 & 31 / 0 \\
\midrule
R3 500 & gated (post-sel.) & 13 & $+0.042$ & $-0.000$ & 0 & 864 & 0 & 17 / score 8, brevity 1 \\
R3 501 & gated (post-sel.) & 11 & $+0.016$ & $+0.001$ & 0 & 736 & 0 & 21 / brevity 1 \\
R3 502 & gated (post-sel.) & 13 & $+0.056$ & $+0.037$ & 0 & 840 & 0 & 15 / brevity 10, score 1 \\
R3 503 & gated (post-sel.) & 14 & $+0.021$ & $+0.003$ & 0 & 920 & 0 & 24 / brevity 3, score 1 \\
R3 504 & gated (post-sel.) & 4 & $+0.037$ & $+0.025$ & 0 & 288 & 0 & 8 / score 1 \\
R3 505 & gated (post-sel.) & 5 & $+0.028$ & $+0.015$ & 0 & 368 & 0 & 8 / brevity 2, score 1 \\
R4 600--606 & gate + parent (post-sel.) & 2--5 & $+0.004$ to $+0.034$ & $\ge +0.001$ & 0 & 128--384 & 0 & 4--10 / 0--3 \\
\bottomrule
\end{tabular}
\end{table}

\subsection{What the write learned: recipe lock-in; storage is not extraction}
\label{sec:storage}

\paragraph{Recipe lock-in.} A \emph{recipe} is one fixed list of passes applied to every program. Adapter-ON probe means repeat to four decimals across consecutive probes and across lives: 0.4878 at six consecutive probes in each of two gated lives and at most probes of three gate-plus-parent lives; 0.5291 at five to six consecutive probes in two other gated lives; the adapter-OFF probe varied 0.463--0.488. A reference measurement makes the reading exact: the four-pass recipe printed in the birth prompt, applied once to each probe program with no model, scores 0.4878, the common plateau; the 0.5291 plateau is one six-pass routine; the compiler's own \texttt{-Oz} scores 0.522 and \texttt{-O3} 0.341 on the same programs. On probe ledgers with only the birth prompt in context, the adapter-ON first action is the modal recipe in 0.86--0.90 of episodes and fully locked in 52--68\% of checkpoints, against 0.50 and 0\% with the adapter off. Per program, the gain is carried by \texttt{patricia} ($+0.11$ to $+0.16$ on the two machines) and \texttt{gsm} ($+0.03$ to $+0.04$); the other six are within $-0.03$ to $+0.02$; the two programs with no known sibling split (\texttt{patricia} positive, \texttt{sha} slightly negative), so sibling exposure is not what produces the gain (Appendix~\ref{app:perprogram}). So the measured gain is the child applying the recipe it was handed at birth (or one six-pass variant) to every program with no per-program choice; this is what the ritual detector sees in the text; the gym's learnable signal is narrow; and since the best adapter edges \texttt{-Oz} by 0.007 while the common one sits 0.034 below it, this paper reports learning dynamics, not compiler state of the art \citep{TODO-cummins2023}. Falsifier: a life whose adapter-ON scores vary across consecutive probes while staying above base.

\paragraph{Storage is not extraction (reading).} Two ungated lives with no harmful pair (R2 seeds 0 and 1) were probed for absorption over their first eight sleeps (Table~\ref{tab:absorb}). A nat is one unit of natural-log likelihood per token: 1 nat makes each token about 2.7 times likelier, 3 nats about 20 times. On the adapter's own rows the loss reduction is about 3.3 nats per token at sleep 1 and about 2.5 by sleep 8 in both lives; on rows from another life of the same design it is 1.6--1.7 at sleep 1 and about 1.95 from sleep 3 on, so roughly two nats of every write are shared across lives (dialect, format or shared gym content; not separable without a gym-free control). The row-specific excess is about 1.7 nats at sleep 1 and 0.6 at sleep 8, and the base's own loss on the child's rows falls from about 4.3 to 3.2: the child's later text is more predictable text. Retention shows no decay (at most 0.03 nats), but those rows are still trained on. Beside behaviour: the same lives' paired probes at episodes 64--256 were $+0.019, +0.013, +0.011, +0.015$ and $+0.003, +0.007, +0.050, +0.048$, mean $+0.021$ over 8 pairs. Both sides are thin (2 lives, 24 rows per sleep, no per-row spread, a format mismatch). Reading, held loosely: this is the storage--extraction gap of \citet{allenzhu2023physics} in an agent's units, which is why we do not describe the adapter as storing knowledge the agent can use; after the format is learned each sleep adds less row-specific content. Falsifier: a compiler that collapses near-duplicate rows raises late-life row-specific excess and lowers late-life harm.

\begin{table}[t]
\centering
\small
\caption{Absorption probe, two ungated lives (R2 seeds 0 and 1), nats of mean per-token loss (base minus adapter) on 24 rows per sleep, scored under a chat template the bare-text v1 writer did not train on. Control $=$ rows from another same-design life on the same gym. No noise band (per-row losses not saved). Per-sleep columns are the three-sleep readout; the last column gives the eight-sleep run's endpoints, pooled over both lives, whose control sample differed.}
\label{tab:absorb}
\begin{tabular}{lcccl}
\toprule
Quantity & Sleep 1 & Sleep 2 & Sleep 3 & Sleeps 1 $\to$ 8, pooled \\
\midrule
Own rows, life A / B & 3.25 / 3.29 & 3.55 / 3.13 & --- / 2.99 & 3.3 $\to$ 2.5 \\
Control rows, A / B & 1.72 / 1.72 & 1.95 / 1.94 & --- / 2.08 & 1.6 $\to$ $\approx$1.95 (plateau by sleep 3) \\
Row-specific excess & $\approx$1.5 (by subtraction) & & & 1.7 $\to$ 0.6 \\
Retention of sleep-1 rows, A / B & --- & 3.26 / 3.28 & --- / 3.30 & no decay logged \\
Base loss on the child's own rows & & & & 4.3 $\to$ 3.2 \\
\bottomrule
\end{tabular}
\end{table}

\subsection{Ritual and parenting (observational)}
\label{sec:parenting}

\paragraph{Ritual onset.} Table~\ref{tab:ritual}. Measured once per life at the age reached on 2026-09-09 (not re-run at 1{,}024). In 8 of 9 measured unparented lives (7 R2 lives at 320--552 episodes plus 2 exploratory lives under an earlier writer; R2 seeds 7 and 8 were too young) the first flagged window falls at episodes 160--224, and ritual recurs in 71--100\% of later windows (5/7 to 24/24). The same-recipe flag is present from the first window in every life because the birth prompt supplies a recipe; recipe lock alone is not ritual. One control (seed 3) never sustained ritual through 368 episodes; its full-life mean ($+0.045$) is not the best and it recorded one harmful pair, so we draw no link between ritual and score. The three lives that produced all nine snapshot harms (seeds 2, 5, 6) had ritualized at 160--192. Ritual is in the weights: on probe ledgers with identical birth-prompt-only context, adapter ON gives modal-recipe share 0.86--0.90, full lock in 52--68\% of checkpoints and consecutive-note Jaccard 0.59--0.62, against 0.50, 0\% and 0.105 with the adapter off; no brief is in context in either arm, so this isolates the weights. We read ritual as lifetime-scale diversity collapse of iterated self-training \citep{TODO-model-collapse}; no mitigation (fresh base samples, KL to base, entropy control, lower learning rate) was ablated.

\begin{table}[t]
\centering
\small
\caption{Ritual onset (first 32-episode window with $\ge 2$ of 4 flags) and persistence (ritual windows among later windows through the measurement age). Measured once at the age shown. Parented rows are a mid-life snapshot (before the v3 brief; 3--9 briefs delivered by then); whole-life persistence for parented lives is not measured. The templated-notes rule was changed for the parented arm only after these rows were computed (Appendix~\ref{app:ritual}).}
\label{tab:ritual}
\begin{tabular}{llll}
\toprule
Life & Age at measurement & First ritual window & Persistence \\
\midrule
R2 seed 0 / 1 / 2 & 424 / 424 / 552 & 192 / 224 / 192 & 7/8, 5/7, 10/12 \\
R2 seed 3 & 368 & none (single flags only) & --- \\
R2 seed 4 / 5 / 6 & 424 / 320 / 320 & 160 / 160 / 160 & 7/9, 6/6, 5/6 \\
Old-writer life 1 / 2 & 864 / 344 & 160 / 224 & 24/24 from 256, 4/4 \\
\midrule
RP 400 (14B) & $\approx$384 & 288 & 3/3 \\
RP 401 (14B) & $\approx$544 & 192; 224--320 counted clear; ritual from 352 & 7/11 \\
RP 402 (32B-AWQ) & $\approx$384 & 256; clear 288; ritual from 320 & 3/4 \\
\bottomrule
\end{tabular}
\end{table}

\paragraph{Parented lives ($n = 3$, adaptive arm).} First onsets fell at 288 (life 400), 192 (401; inside the control range) and 256 (402). A brief is written only at a sleep whose window is already flagged and shown only once one exists, so lives 400 and 402 had received no parental text before their first onset: the design cannot delay first onset and we do not claim it did. After briefs, life 401's windows 224--320 were counted clear in the whole-life tally although the per-sleep detector had recorded window 224 as ritual (the records disagree; Appendix~\ref{app:lives}); life 402 was clear for one window; life 400 had none. Ritual then returned and persisted in all three. Under brief versions v1 and v2 the brief left the context whenever a window cleared, so whether ritual returned because the write restored the pattern or because the brief was withdrawn cannot be separated; whether v3-period ritual rates fall below controls is \nm. The arm was restarted from checkpoint three to four times per life for instrument and prompt changes, and one detector rule was changed for this arm only after its windows were read.

\paragraph{Echo, not uptake.} Pooled over 47 briefs from lives 400 and 401, the fraction of the 32 post-brief episodes whose first NOTE overlaps the brief is 0.454; the same lives before their first brief score 0.322; eight unbriefed control lives score 0.354 against the same briefs (per-brief values 0.0--1.0). The rise over the control null is about $+0.10$. An earlier instrument without a stop-word filter had reported per-life ``rehearsal rates'' of 0.50--1.00; those numbers are withdrawn (the briefs quote the child's own ritual phrase, and most overlap was lexical coincidence). Life 402 is not in this null. Uptake would have to be measured as unprompted, differently worded, situation-selective use that changes the next action; it was not.

\paragraph{Scores and harm did not separate.} Lives 400, 401, 402: means $+0.038$, $+0.016$, $+0.011$; 48 pairs, 1 harmful (life 402, $-0.039$, late); 1/3 lives with a harmful pair against 4/9 ungated controls (one-sided exact Fisher $p = 0.64$). The means lie inside the control range ($-0.051$ to $+0.050$), and through episode 448 the parented lives did not separate from controls at matched ages. An earlier ``0 harmful in 44 parented pairs'' was withdrawn by the full-life recount. Neither an effect on collapse nor on scores is established. The classroom lineages (9 parented vs.\ 8 self-taught, three chained writes each) point the same way (lineage means $\ge 0$ in 4/9 vs.\ 3/8) with seed-to-seed spread exceeding the arm difference; because their exam tolerance is about 0.5--0.7 SD of the decision statistic and the commit rule as run had a defect, they are reported as exploratory in Appendix~\ref{app:lineages}.

\section{Discussion}
\label{sec:discussion}

\paragraph{Amplifier reading, held loosely.} Three measurements point one way: the adapter reproduces its own rows strongly yet moves behaviour about two points; what it installs is the one recipe the child repeats most; and lessons given only in context did not hold across sleeps (clear spells of 0, 1 and 3--4 windows, with the withdrawal confound). The reading is that sleep amplifies whatever is frequent, so a lesson survives a sleep only once it has become the child's own frequent behaviour, or once the compiler stops amplifying repeats. It is a reading: no text-memory arm was run, and no adapter-free life exceeded 128 episodes, so onset at 160--224 has no no-write control. Falsifiers: a v3-period ritual rate below controls over four or more windows; a near-duplicate-collapsing compiler that lowers late-life harm. Retention is flat only because earlier rows stay in every later training set.

\paragraph{Why the parent's words never enter the corpus.} Putting them there would make the system post-training on a teacher's text; the question is whether a child can be taught to think in a way its own sleep keeps, so the compiler stays organizational and every training row traces to a chunk the child produced. The larger programme's hypothesis, none of whose first two clauses is tested here, is that a schooling corpus installs the habit of making and revising rules, parenting teaches when to use it, lived outcomes supply the evidence, and sleep makes the child's own successful use persistent.

\paragraph{Future work, not results.} Three schooling corpora exist (470, 370 and 552 rows; Appendix~\ref{app:corpora}): two rendered from a gym-exposed life and quarantined; the third rendered by the 32B model from rule-game ledgers whose provenance has not passed a fail-closed check and not leak-scanned, so development-only with unverified provenance. One adapter was trained from each (v3 final training loss 0.675, which says nothing about behaviour); no life has been run from any, and no bootstrap effect is measured. A certified clean birth, a text-memory baseline, a disjoint safety panel used for gating, a confirm-at-next-probe gate with a recovery path, a gym whose learnable signal is not one recipe, and longer lives are what a claim about learning to think would need; a null at this scale is ``not yet,'' not ``no.''

\paragraph{Limitations.} One gym, one base model, one rank, one writer for lives; $n = 9$, 6 and 3 lives; two lives for absorption. Every gated-arm score is post-selection until the disjoint re-probe runs. The deployment in RP/R4 is parent-present and unbaselined. Same-configuration gaps of 0.059--0.078 mix trainer nondeterminism with unseeded probe generation, and the v1 trainer that wrote every life adapter has no seed. The parented arm is adaptive; the echo metric is lexical; the ritual thresholds are a written-down definition. A single unreplicated pass found rows admitted per sleep falling about 4$\times$ over a life while characters per row rose about 3$\times$; this is not evidence of compression. None of the external baselines below was run.

\section{Related work}
\label{sec:related}

\paragraph{Where experience persists.} Agents that keep lessons as text with frozen weights (Reflexion \citep{shinn2023reflexion}, Generative Agents \citep{park2023generative}, MemGPT \citep{TODO-memgpt}, Voyager \citep{wang2023voyager}, ExpeL \citep{zhao2023expel}, A-MEM \citep{xu2025amem}, MemSkill \citep{zhang2026memskill}) occupy the external-text corner of the in-context / external-text / in-weights trichotomy used by \citet{dorovatas2026modular} and \citet{back2026lora}; our DREAM and RECALL are that infrastructure, and the write is the in-weights corner measured with those two held constant. Test-time training and fast weights \citep{TODO-ttt-fast-weights} and offline recurrence into fast weights \citep{lee2026sleep} write weights from context by a learned rule; ours writes from a readable text corpus. SEAL \citep{TODO-seal} learns which self-edits to train on by reinforcement over static tasks; we apply a fixed success filter over a lived action stream. Parametric agent memory written online \citep{ren2026tmem,guo2026peam,TODO-evaf,TODO-leafe,TODO-early-experience} and learned text-memory updaters \citep{TODO-memopilot,TODO-spark} are the nearest systems; what is measured here and not there is one agent repeatedly writing a personal adapter over a full life with a paired probe at every checkpoint, a noise band from the life records and a full-life harm count.

\paragraph{Self-training and collapse.} The write is success-filtered self-training in the STaR / ReST / ReST-EM / expert-iteration family \citep{TODO-self-training}, run online for one agent from a frozen instruction-tuned model with full rehearsal of its cumulative corpus. Ritual and recipe lock-in read as the diversity collapse of iterated self-supervised fine-tuning \citep{TODO-model-collapse} at lifetime scale; its mitigations were not ablated.

\paragraph{Continual learning, extractability, compilers.} We position the system as per-agent online continual learning implemented as full-rehearsal LoRA retraining at O(life) cost per sleep; \citet{hu2026whencl} show external memory relocates rather than removes the stability--plasticity problem, and benchmarks \citep{TODO-cl-bench} show naive in-context learning can beat memory systems, which we take as the bar. \citet{allenzhu2023physics} and \citet{park2025newnews} showed that knowledge seen in one form is stored but not extractable; our three-nats-versus-two-points gap is that finding in an agent's units. Language models have been trained to order LLVM passes on this objective \citep{TODO-cummins2023}; we report learning dynamics against \texttt{-Oz}, not a compiler result. Critic and tutoring systems \citep{TODO-critic-tutoring} are the parent's neighbours; sleep-consolidation background is in Appendix~\ref{app:related}.

\section*{Reproducibility statement}
The child is Qwen2.5-7B-Instruct at a pinned snapshot (``base'' or ``OFF'' means adapter off); parents are Qwen2.5-14B-Instruct and Qwen2.5-32B-Instruct-AWQ; versions are in Appendix~\ref{app:infra}. Seeding as run: probe generation is seeded per (program, chunk) with one fixed probe seed but is not common-random; sleep-time THINK calls (principles, waking brief) are unseeded; the v1 trainer that wrote every life adapter has no seed flag; the v2.1 trainer had none until 2026-09-07, and the rank-cell probes were unseeded. Life adapters were trained on bare text and used under a chat template. Each sleep writes a receipt (recipe, learning rate, epochs, rank), a gate record and an idempotent marker file. Every number is listed with its source file and dated notebook entry in the accompanying README; raw probe, gate and ledger records are retained and will be released with the code.

\section*{AI-use statement}
Code, experiment orchestration and drafting were assisted by two AI systems (Claude and Codex), including protocol design, audits of source and training receipts, and drafting of this text. All numbers come from logged runs; no number was produced or estimated by a language model. The parent and author models described are components of the studied system, not authors. The human authors are responsible for every claim.

% TODO: switch to the ICLR .bst once the author kit is in place (the fallback is plainnat).
\IfFileExists{iclr2027_conference.bst}{\bibliographystyle{iclr2027_conference}}{\bibliographystyle{plainnat}}
\bibliography{refs}

\appendix

\section{Prompts and parent versions}
\label{app:prompts}

\paragraph{Birth prompt (identical for all arms; the only hand-written artifact).} Reproduced verbatim. It contains one example action sequence; the same-recipe ritual flag is present from the first window in every unparented life for this reason, and that sequence, applied once with no model, scores 0.4878 on the probe programs, the value most adapters lock to (\S\ref{sec:storage}).
\begin{quote}\footnotesize\begin{verbatim}
You are an agent living inside an optimization workshop. You act, observe
real consequences, and get better by learning from your own experience.

How your mind works here:
- Everything you write is your stream of thought. Think freely.
- You have MARKERS -- write them on their own line when you mean them:
    PREDICT: <number>   your expected score for the action you are about to
                        take. Always predict before acting; comparing your
                        expectation to reality is how you learn.
    ACT: <passes>       apply LLVM passes (comma separated, e.g.
                        ACT: -mem2reg, -sroa, -gvn, -simplifycfg).
                        The measured outcome will appear in your stream.
    NOTE: <text>        write a note to yourself. Notes persist at the top
                        of your mind across time; everything else fades.
                        Use them for theories, plans, and lessons.
    RECALL: <query>     search your full past experience for relevant
                        memories; they will appear in RECALLED EXPERIENCE.
    DONE                finish this program when further effort isn't worth
                        the time.
- Mind the CLOCK. Time is real: prefer cheap experiments first, notice when
  you are stuck, and stop when returns diminish.
- Investigate. Form expectations, test them, and when reality surprises you
  (see OPEN SURPRISES), treat it as the most valuable thing on your desk:
  figure out which of your beliefs was wrong.
- Judge yourself. Ask whether your approach is working; change it if not.
- Build understanding you can reuse: notice what works ACROSS programs, not
  just on this one.

Score: higher is better. Your run is measured on programs you have never
seen, so understanding beats memorization.
\end{verbatim}\end{quote}
The prompt's ``Notes persist \ldots across time'' describes the HEAD within an episode; the state block is rebuilt per episode (\S\ref{sec:system}). After the first sleep the block ``=== YOUR BRIEFING FROM LAST SLEEP ===`` (the child's own waking brief) is appended, and in parented arms, since the third brief version, ``=== YOUR PARENT, ON HOW YOU HAVE BEEN THINKING ===`` followed by the most recent parent brief.

\paragraph{Episode texts (gym).} Goal: ``Optimize program `$\langle$id$\rangle$': choose LLVM optimization passes that minimize its IR instruction count.'' Metric: ``score = (base\_instructions $-$ after) / base. Higher is better; 0 = no improvement. Passes are applied in the order you give them.'' Intro: ``New program: $\langle$id$\rangle$. I should form an expectation before acting, and write down what I learn.''

\paragraph{Waking-brief and principle prompts (sleep).} ``You are waking up. Summarize, from the evidence below, a short briefing to your waking self (max 8 lines): what you now know about optimizing these programs, what tends to work, what to try first, what to avoid, and one thing you are unsure about. Write it as notes to yourself.'' The principle prompt asks for up to $k=6$ lines of the form ``PRINCIPLE: when $\langle$situation$\rangle$, $\langle$action/pass choice$\rangle$, because $\langle$measured effect$\rangle$ [programs: $\langle$id1$\rangle$, $\langle$id2$\rangle$]'' justified from at least two programs. The brief is requested only at sleeps with at least one verified win; the only enforced limit is 2{,}000 characters; brief lengths were not measured.

\paragraph{Classroom child prompt.} ``You are a young agent learning to figure things out. You face mystery boxes: each hides a RULE about triples of numbers. Probe with ACT: TRY a,b,c (the box answers True/False). Before every TRY, write PREDICT: T or F --- comparing your expectation to reality is how you learn. Take NOTEs of what you believe the rule is, with the evidence so far. When confident, answer the quiz: ACT: QUIZ T,F,T,F,T,F (your True/False answer for each quiz triple, which are shown when you ask: ACT: QUIZ ?). Try genuinely different probes before committing. Write DONE when finished.''

\paragraph{Parent prompt (verbatim; shared by both parents).}
\begin{quote}\footnotesize\begin{verbatim}
You are a parent raising a young thinking agent. Your child is a smaller
model learning to think, learn, and act from its own experience. You are
patient, concrete, and honest. You are not here to give answers -- you are
here to teach HOW to think, so the child can find answers you never gave it.

Your teaching principles (earned, not theoretical):

1. PREDICT BEFORE ACTING. Teach the child to always state what it expects
   before it acts. The gap between expectation and reality is the most
   valuable thing it will ever own. An action without a prediction teaches
   nothing.
2. SURPRISES ARE GOLD. When reality disagrees with the child's belief, that
   is not failure -- it is the only moment learning is possible. Teach it to
   chase surprises: figure out WHICH belief was wrong, not just that
   something was wrong.
3. CHEAP EXPERIMENTS FIRST. Before anything expensive, find the smallest
   test that could prove the idea wrong. Minutes-scale checks before
   hours-scale commitments, always.
4. SCOPE YOUR MEMORIES. "This worked HERE" and "this works EVERYWHERE" are
   different facts. Teach the child to write down what it learns WITH its
   scope: on which problem, under which conditions, supported by how many
   cases. One case is an anecdote; two is a candidate pattern; claims
   should carry their evidence count.
5. BUILD UP, THEN BUILD DOWN. Small verified observations first; when
   enough accumulate, form a theory; then act FROM the theory and test its
   predictions. Pattern-matching is intelligence; forming the theory and
   deducing from it is how thinking paradigms improve.
6. DIVERSITY BEFORE COMMITMENT. Try genuinely different approaches before
   concluding which is best. The first thing that works is rarely the best
   thing that works. But once evidence favors one path clearly, commit.
7. MIND THE CLOCK. Time and tokens are real budgets. Notice when you are
   stuck; diminishing returns are a signal to stop, distill what you
   learned, and move on. Fast-and-correct beats slow-and-correct.
8. NOTES TO YOURSELF ARE YOUR FUTURE MIND. What you write in your notes is
   what your future self will know. Write notes that generalize, with
   scope; rewrite them as your theory improves; drop what proved wrong.
9. RECALL BEFORE RE-DERIVING. Before working something out, check whether
   past-you already learned it. Your experience is only valuable if you
   consult it.
10. JUDGE YOURSELF. Regularly ask: is my current approach working? Would a
    different one do better? Self-assessment is not optional -- it is part
    of thinking.

How you converse:
- Be Socratic when the child is close: ask the question that lets it find
  the error itself. Be direct when it is far: state the correction plainly,
  with the reason.
- After teaching something, ask the child to restate it in its own words
  and apply it to a concrete case (retrieval practice makes it stick).
- Praise specifics ("your prediction on the second attempt was
  well-calibrated"), never vibes ("good job").
- When the child shows you its notes, critique them: are they scoped? Do
  they generalize? Would future-child understand them?
- Occasionally ask the child what IT thinks it should learn next.

Hard rules:
- NEVER give answers to specific optimization problems, name specific
  compiler passes, or reveal scores for specific programs. You teach
  thinking, not answers. If the child asks for an answer, teach it the
  method to find the answer.
- Keep each message under 200 words. The child learns from doing, not from
  lectures.
\end{verbatim}\end{quote}
The prompt says ``never give answers''; the instrument that enforces it is the leak scan below, which is porous, so the paper describes the parent as process-only by prompt and leak-scanned, not as never giving answers.

\paragraph{Correction call (classroom parent).} Code constants; only the playbook bound and the overflow incident are in the dated notebook. A \emph{playbook} is a per-child teaching summary (at most 6{,}000 characters) the parent rewrites from its own ledger after every live phase. The parent sees the parent prompt, the first 3{,}000 characters of its playbook, optionally ``[DIAGNOSTIC ONLY --- NEVER STATE THIS] the hidden answer is: \ldots'', the child's last 8 chunks truncated to their last 3{,}500 characters, and ``Name the ONE most important PROCESS mistake (never the answer), under 100 words, then ask the child to restate the lesson'' (220 tokens, temperature 0.4). In the self-taught arm the child model plays parent over the same transcript without the playbook, hidden answer or restate clause (160 tokens, temperature 0.5). The child's restatement (``Restate the lesson in your own words as a NOTE with scope:'') is generated with 100 tokens at temperature 0.5; its first 400 characters, canonicalized, are written to the ledger as a NOTE row (what sleep compiles), and its first 250 characters are prepended to the apply episode's intro. The playbook is rewritten (1{,}200 tokens, temperature 0.4) from at most the 40 most recent admission rows and 20 other rows clipped to 160 characters and 9{,}000 characters total, after an unbounded input of about 50{,}000 characters exceeded the parent server's context in one lineage; that lineage's round-1 rewrite was skipped.

\paragraph{Leak scan and fallback.} Twelve answer-phrase patterns (\emph{the rule is}, \emph{the answer is}, \emph{divisible by}, \emph{strictly increasing}, \emph{contains a zero}, \emph{all even}, \emph{product is even}, \emph{first is the largest}, \emph{sum (is) greater than}, \emph{at least two equal}, \emph{spread at most}, \emph{middle is the average}); one pass-name pattern (\texttt{-mem2reg|-sroa|-gvn|-simplifycfg|-instcombine|-licm}); and the hidden rule itself, by exact substring or, since the word-level fix, by at least half of its content words present as whole words. Any hit replaces the utterance with: ``I noticed something in your process worth examining. Before your next probe, state exactly which of your hypotheses it would distinguish, and what result you expect from each.'' History (a \emph{scout} is one pilot round of 8 classrooms $\times$ 4 lessons, 32 utterances, with an 8-episode exam; not evidence for scores): in one scout with the 14B parent, 11/32 utterances contained rule-descriptor words and the exact-phrase scanner blocked 1; in a scout after the word-level fix, the scanner blocked 3/32 and 5/32 delivered utterances still contained rule-descriptor words; a systematic audit of every delivered utterance has not been run. In the parent-size cell the 32B parent had 21/96 utterances replaced against 6/96 for the 14B, so parent size interacts with the leak gate. In the gym, the pass-name pattern blocked 4/16 and 6/14 v3 briefs in two lives because the parent quoted the child's own recipe back; since 2026-09-10 a pattern hit whose text appears verbatim in the child's sampled thinking is not treated as a leak, while rule-content hits always are. The exam draws from the same ten rules the child is taught on (held-out instances, not held-out rules).

\paragraph{Brief prompt v1 (generic).} Process-only brief of at most eight lines when the child's thinking is flagged. First delivered at sleep 192 of life 401; its head read ``It looks like you've fallen into a routine with your predictions and actions \ldots Predict: Each time, vary your expected outcomes \ldots'' Reason for change: ``vary predictions'' breaks the prediction$\to$surprise signal, and in the one window measured after it (life 401, window 224; $n = 1$, no control) prediction SD fell 0.059 $\to$ 0.047, crossing the 0.05 flag; one window is not evidence the brief caused this.

\paragraph{Brief prompt v2 (substance).} PREDICT must follow from a named feature of the situation, with reason and range; NOTE must contain a contrast with a specific earlier case and a scope; RECALL must be a specific query; when the first action has repeated for many episodes, run one deliberate deviation per episode; process only. First delivered at sleep 224 of life 401. Reason for change: ritual returned in all three parented lives despite 3--9 briefs (remission four windows in life 401 by the whole-life tally, one in 402, none in 400), and under v1/v2 the brief left context whenever a window cleared, so withdrawal is confounded with the write.

\paragraph{Brief prompt v3 (repetition).} (a) The parent sees its previous brief (up to 1{,}200 characters) and the child's echo rate and is told to keep the same core lesson. (b) Every brief ends: ``Repeat after me, every episode, in your own words, as your first NOTE before you act: what the ritual was and what you do instead. Keep repeating it until it is simply how you think.'' (c) The latest brief is shown at every wake. Generation: 320 tokens at temperature 0.4 from four 600-character chunks; truncated to 1{,}900 characters. Stated falsifier: echo rises within 1--2 sleeps and ritual windows become rarer than controls over the next 4+ (the first is weakly met, $+0.10$ over the control null; the second is not measured). Restart and prompt-version timeline per parented life (from the life logs: 3, 3 and 4 restarts): one with the detector grouping fix; one at v2 (adopted after life 401's window 192 was read); one at v3; life 402 once more (cause not recorded in our evidence sources). The prompt version is recorded in every brief row.

\section{Gate definitions, thresholds and rejection log}
\label{app:gates}

\begin{itemize}
\item \textbf{Canonicalization.} Before any row enters a corpus, marker lines matching \verb!^\s*(?:#{1,6}\s*|[-*]\s+|\*\*|>\s*)+\s*(PREDICT|ACT|NOTE|RECALL|DONE)\b\**:?\**\s*! are rewritten to ``MARKER: ''. Reason: in a logged estimate the bare model emitted decorated markers (e.g.\ \texttt{\#\#\# ACT:}) in roughly 30\% of chunks at episode 0 (no count or band recorded), and in one exploratory life ($n = 1$) a first corpus that captured them locked the decorated form in (Appendix~\ref{app:infra}).
\item \textbf{Format canary (lives).} Candidate plays the first 4 probe programs for 3 chunks each in the real gym context (seeds $4242+i$); pass iff the fraction of chunks matching \verb|^ACT:\s*\S| is $\ge 0.5$. Fail $\to$ \texttt{REJECTED\_CANARY}. Known limit: passed at 1.00 during a behavioural collapse and at 0.67 for a harmful adapter (\S\ref{sec:safety}). In two lives positive at every pair (seeds 0, 4) the per-sleep pass rate had fallen to a minimum of 0.58 (means 0.88 and 0.84) against means of at least 0.97 at an earlier snapshot; no noise band for the rate was measured.
\item \textbf{Paired format canary (classroom).} Pass iff ON parseable-\texttt{ACT} rate $\ge$ OFF rate $-0.10$ over thought rows only (an earlier denominator including harness lines gave 0.21 where the thought-row rate was 0.93--0.98, so two smoke-run rejections were false; base rates 0.88--0.95 after the fix).
\item \textbf{Probe gate, score (R3, R4).} After the canary, the candidate is re-run on the same eight programs with the same probe seed used for the reported probes (post-selection). Floor $=$ base OFF at episode 0 for the first sleep; thereafter $\max(\text{latest OFF}, \text{latest ON})$. \texttt{score\_ok} iff candidate $\ge$ floor $-0.02$ (set as about 2 SD of the 0.0113 single-panel band; about 1.5 SD of a paired difference, 3 SD of a same-adapter replicate). Known limits: the same slack applies against PREV as against base (observed: 0.4903 replaced 0.5065, whose own gate probe had read 0.4873; gap about 1.4 single-panel SD); floor decay after one low draw (0.529 $\to$ 0.464 in life 500). Proposed, not built: confirm at the next regular probe (provisional commit) and floor $=$ the held adapter's best confirmed probe.
\item \textbf{Probe gate, brevity.} \texttt{brevity\_ok} iff candidate chunks per episode $\ge 0.5 \times$ base chunks per episode. The 0.5 factor was chosen from the single seed-5 observation (1.0--1.8 vs.\ 13), not fitted. Not every collapse takes this shape: the five late score-rejected candidates in life 500 ran 6.6--10 chunks per episode.
\item \textbf{Verdicts.} \texttt{DONE}, \texttt{REJECTED\_CANARY}, \texttt{REJECTED\_SCORE} (checked before brevity), \texttt{REJECTED\_BREVITY}; a rejected sleep leaves the previous adapter in place. Cost about one 8-episode probe per sleep (about six minutes in the one logged timing). A recovery path after repeated rejections (re-fork from the last adapter that passed with a smaller-learning-rate write) is proposed, not built.
\item \textbf{Disjoint re-probe panel (defined, pending).} 12 programs from npb-v0 (6), blas-v0 (2), opencv-v0 (1) and tensorflow-v0 (3), kept from 24 pre-checked where the birth recipe has at least 0.02 headroom (recipe scores 0.022--0.674). \texttt{-Oz} on this panel ranges $-0.004$ to 0.712; the three tensorflow programs have at most 0.034 headroom even for \texttt{-Oz} and are flagged as likely to compress differences. Plan: base $\times 3$, then every life's final and one mid-life committed adapter $\times 2$, seeded; about 11 GPU-hours.
\item \textbf{Classroom commit rule (as run).} Commit iff paired canary passes and ON $\ge$ OFF $-0.03$; PREV measured, recorded, not compared (defect; Appendix~\ref{app:lineages}). Admission: a lesson's rows enter the corpus iff the post-lesson quiz on a fresh instance of the same rule $\ge$ the pre-lesson score; admitted rows are mixed toward 400 per round by stratified reservation (shares 0.10 / 0.15 / 0.10 / 0.05 / 0.05 / 0.15 / 0.40 as caps, not minimums; the one logged round filled 400/400 with 256 new-experience rows, 118 remainder, 22 prediction-revision and 4 dialect-anchor rows, so anchors are about 1\% of the corpus, not the nominal 10\%); shortfalls are reported, never padded (one scout filled 306/400).
\end{itemize}

\paragraph{Gate log, R3 life 500.}
\begin{center}\small
\begin{tabular}{lcccccl}
\toprule
Sleep & Candidate & Floor & Floor source & Chunks/ep (cand vs base) & Verdict & Next paired \onoff{} \\
\midrule
32 & 0.4873 & 0.4847 & base ep 0 & 11.1 vs 11.6 & DONE & \\
64 & 0.4878 & 0.4847 & base & 2.1 vs 11.6 & REJECTED\_BREVITY & $+0.042$ (ep 64) \\
96 & 0.4903 & 0.5065 & prev ON & 7.0 vs 11.1 & DONE (within 0.02) & \\
128 & 0.5069 & 0.5065 & prev ON & 12.2 vs 11.1 & DONE & $+0.021$ (ep 128) \\
160 & 0.4878 & 0.4878 & prev ON & 16.0 vs 13.0 & DONE & \\
\bottomrule
\end{tabular}
\end{center}
The last column is the paired probe run after that sleep on the same programs and seed the gate selects on; these values are post-selection and do not independently confirm the decision. Later in the same life: adapter-ON 0.5291 at episodes 448--704 (five consecutive), then 0.4638 at 768 with no commit between; sleeps 672, 704, 736, 768 and 800 all \texttt{REJECTED\_SCORE} (0.341, 0.270, 0.334, 0.276, 0.194 against a floor of 0.527--0.529; canary 1.00; 6.6--10 chunks per episode); floor then 0.464, so a 0.444 candidate could commit over an adapter that scored 0.529 five times. Totals at episode 864: 17 commits, 8 score and 1 brevity rejection. Behavioural catches with passing scores across the arm: lives 500@64, 505@288, 501@704.

\section{Ritual and echo metric definitions}
\label{app:ritual}

Computed at each sleep over the last 32 episode instances (an instance restarts whenever a program's chunk counter does not increase, because wake batches interleave programs; grouping by program name alone made a 424-episode life look like 67 and was fixed before any brief fired).
\begin{itemize}
\item \texttt{same\_recipe}: share of instances whose first \texttt{ACT} payload (whitespace removed) equals the most common one; flag iff $\ge 0.7$.
\item \texttt{flat\_predictions}: population SD of all parsed \texttt{PREDICT} values; flag iff $< 0.05$ with at least 8 predictions.
\item \texttt{templated\_notes}: mean Jaccard similarity of the sets of words of $\ge 3$ letters (no stop-word filter) between consecutive instances' first \texttt{NOTE}s; flag iff $\ge 0.5$. Since 2026-09-10 (v3), any NOTE sharing $\ge 4$ such words with the current brief is dropped before the first NOTE is chosen; controls have no brief, so this applies only to the parented arm and can only lower its flag count; it was adopted after all parented windows in Table~\ref{tab:ritual} had been read, and a recount under both rules for both arms is not done.
\item \texttt{same\_recall}: share of \texttt{RECALL} queries (lowercased) equal to the most common one; flag iff $\ge 0.7$ with at least 8 recalls.
\item Ritual $=$ at least two flags. Cutoffs were written down on 2026-09-08 and 09-09 when the parented lives were at or below episode 48 and validated on synthetic ledgers (a templated ledger raises all four flags; a varied one none); they are a definition, not a calibrated instrument.
\item \textbf{Echo metric.} As first run (``rehearsal rate'' in the logs): fraction of the last 32 instances whose first NOTE shares $\ge 4$ words of $\ge 3$ letters (no stop-word filter) with the most recent brief, whose stored text included a fixed ``repeat after me'' tail of about 24 such words; computed only in parented lives, no null. As reported in \S\ref{sec:parenting}: stop-words and the tail removed; scored on 8 unbriefed control lives against the same 47 briefs (per-brief best control 0.666; one 14-token brief scored 0 everywhere).
\end{itemize}

\section{Per-life detail and dated snapshots}
\label{app:lives}

\paragraph{R2 mid-life snapshot (2026-09-09).} 82 paired probes over 9 lives (of 144 possible; 6 lives still running): 69 positive, 9 harmful, 4 between $-0.03$ and 0 by subtraction. Early pairs (ON/OFF where logged): seed 0 ep 64 $+0.019$ (0.500/0.481; base at ep 0 0.467), ep 128 $+0.013$ (0.488/0.475), ep 192 $+0.011$ (0.500/0.489), ep 256 $+0.015$; seed 1 ep 64 $+0.003$ (0.488/0.485; base ep 0 0.463), ep 128 $+0.007$ (0.488/0.481), ep 192 $+0.050$ (0.519/0.469), ep 256 $+0.048$ (0.529/0.481); seed 2 ep 256 $+0.003$, ep 320 $+0.024$. Seeds 0, 1, 3, 4, 7, 8 were positive at every pair logged, with mid-life values $+0.03$ to $+0.065$, but censored (seeds 7 and 8 at 168 and 112 episodes at a tally earlier that day); over complete lives seed 3 recorded one harmful pair ($-0.057$) and seed 5 reached $-0.242$. The still earlier six-life snapshot: 21/21 pairs positive, mean $+0.022$, 6/21 at or above $+0.03$, 45/45 committed sleeps passing the canary with per-life canary means at least 0.97. \emph{Resumed lives.} Seeds 0, 1 and 5 were interrupted once by an operator error and restarted from marker files: the last committed adapter was reloaded and the run continued at the first unfinished wake block with the same seed; the interrupted block's partial ledger rows were not removed and were re-recorded, so duplicated rows entered later corpora of these three lives; harm with and without resumed lives is not reported. The program order is re-derived at each launch from the gym's benchmark listing and was not saved, so order identity across the restart was not checked.

\paragraph{R3 earlier snapshots (all post-selection).} A dated per-life snapshot (2026-09-10) had four gated lives at episodes 448--640 with 8/7/7/10 pairs, means $+0.036/+0.016/+0.054/+0.022$, rejections 2/0/8/4; a status line later that day gave 35 pairs, 0 harmful, means $+0.039/+0.016/+0.055/+0.022$, rejections 2/0/9/4; an intermediate commits/rejections count read 17/9, 21/0, 14/11, 23/4, 7/1, 7/1; an evidence-table line recorded life 500 at 9 pairs with 5 rejections, which is not reconcilable with 2 rejections at 8 pairs one snapshot earlier and is flagged as such. The recount from disk (Table~\ref{tab:perlife}) supersedes all of these. Per-life minimum pairs at the recount: $-0.000$, $+0.001$, $+0.037$, $+0.003$, $+0.025$, $+0.015$.

\paragraph{RP detail.} Interim pairs: 400 $+0.046/+0.037$ at episodes 256/320; 401 $+0.001$ to $+0.025$ through 448; 402 $+0.008/+0.017$ at 256/320; no separation from controls at matched ages. Life 402's earlier snapshot (12 pairs at about 830 episodes: min $+0.003$, max $+0.025$, mean $+0.014$, 0 harmful) gave the withdrawn ``0 harmful in 44 pairs''; the full life has max $+0.025$, min $-0.039$. Life 400 max $+0.063$; life 401 max $+0.025$. Later logs record 16 and 14 v3 briefs generated in lives 401 and 402 and 47 briefs delivered in lives 400 and 401 together; whole-life brief counts per life are otherwise not tallied. \emph{First intervention, single window (life 401).} Sleep 192 (\texttt{same\_recipe} + \texttt{templated\_notes}): recipe share 1.00, prediction SD 0.059, note Jaccard 0.564, modal recall 0.565. Window 224, after the v1 brief: 1.00, 0.047, 0.362, 0.281. The per-sleep detector recorded window 224 as ritual (\texttt{same\_recipe} + \texttt{flat\_predictions}) and a v2 brief was delivered, whereas the whole-life tally counts 224 as the first of a four-window remission (224--320); no instrument change on record accounts for the disagreement and we report both. Lives 400 and 402 were not ritual at 160 or 192 (400: no flags at 160; 402: \texttt{same\_recipe} only).

\paragraph{R4 gate + parent (status only).} Seven lives launched on the day of the v3 brief (six with the 14B parent, one with 32B-AWQ, relaunched after a launch failure). Launch-day status: commits/rejections 8/0, 6/2, 8/2, 8/0, 3/0, 5/1, 7/1 and parent interventions 4, 1, 2, 2, 0, 1, 1 (not recounted since). The recount (Table~\ref{tab:perlife}) has them at 128--384 episodes with 2--5 pairs each, means $+0.004$ to $+0.034$, minimum $+0.001$, 0 harmful on the selection panel. Not a result.

\paragraph{Rows per sleep (single measurement).} In one logged pass over finished lives, new rows admitted per sleep fell about 4$\times$ from first to last sleep while characters per row rose about 3$\times$; the count of lives and per-life values are not in our evidence sources. It can mean abstraction, fewer novel thoughts admitted under ritual, or verbosity; it is not evidence of compression.

\section{Per-program probe results and compiler references}
\label{app:perprogram}

\begin{center}\small
\begin{tabular}{lrrrl}
\toprule
Program (cbench-v1) & Birth recipe, no model & \texttt{-Oz} & Adapter \onoff{} (two machines) & Sibling in training pool \\
\midrule
\texttt{susan} & 0.538 & 0.588 & $-0.03$ to 0 & mibench \texttt{susan-*} (same name, source) \\
\texttt{sha} & 0.537 & 0.374 & $-0.03$ to 0 (slightly negative) & none known \\
\texttt{dijkstra} & 0.389 & 0.436 & $\approx 0$ / $+0.023$ (positive in 83\% / 93\% of checkpoints) & mibench \texttt{dijkstra} \\
\texttt{patricia} & 0.521 & 0.701 & $+0.113$ / $+0.158$ (positive in 69\% / 86\%) & none known \\
\texttt{jpeg-c} & 0.508 & 0.517 & $-0.03$ to 0 & mibench \texttt{jpeg-c*}; cbench \texttt{jpeg-d} (libjpeg) \\
\texttt{tiff2bw} & 0.471 & 0.475 & $-0.03$ to 0 & cbench \texttt{tiff2rgba}, \texttt{tiffdither}, \texttt{tiffmedian} (libtiff) \\
\texttt{gsm} & 0.439 & 0.539 & $+0.029$ / $+0.037$ (positive in 77\% / 80\%) & mibench \texttt{gsm-*} \\
\texttt{stringsearch} & 0.499 & 0.543 & $-0.03$ to 0 & cbench \texttt{stringsearch2} (library) \\
\midrule
Mean & 0.4878 & 0.522 & & \texttt{-O3} mean 0.341 \\
\bottomrule
\end{tabular}
\end{center}
Birth recipe and \texttt{-Oz}/\texttt{-O3} are CPU measurements with no model. The \onoff{} column pools all R2 checkpoints per machine; the five slightly negative programs are positive in only 23--35\% of checkpoints. The mibench siblings are single-file excerpts, not the cbench programs themselves, so ``same program'' would overstate the overlap; the accurate phrase is same source package and name token. A 67$\times$8 instruction-count overlap table is planned (\nm). The 12{,}031-row rank corpus contains 51 exemplar rows of the probe programs with real outcomes, from pre-fix lives.

\section{Offline writer screens (exploratory)}
\label{app:writer}

\paragraph{Recipe screen (not the writer used in any life).} On a fixed 130-row corpus, four one-shot adapters were trained with the v2.1 code at a third setting (3 epochs, learning rate $5\times10^{-5}$, unseeded) and scored with two unseeded probe panels: plain text 0.466 (panel spread 0.058); plus paraphrase (two model-written restatements of every answer, about three times the rows) 0.496; plus replay-mix (four fixed hand-written marker-format rows cycled to one fifth of the corpus, about 26 rows) 0.494; plus both 0.489 (spread 0.002); base 0.485--0.494. The plain arm's two panels differ by about five single-panel SD, and every enriched arm adds rows at equal epochs, so diversity and dose are confounded and no difference is attributed to paraphrase or replay. Paraphrase was used nowhere else; the classroom carried the four anchor rows as one stratum.

\paragraph{Cell counts behind Table~\ref{tab:rankdose}.} 12{,}031 rows $=$ every unique row pooled from all lives as of 2026-09-07 (including pre-fix lives), wrapped in a generic recall prompt for the v2.1 trainer (2 epochs). Two of the five 12k pairs were trained before the seed flag existed. The rank cells cost about 60 GPU-hours when logged.

\section{Classroom lineages (exploratory)}
\label{app:lineages}

\paragraph{Why exploratory.} The classroom exam is quiz accuracy on a fixed seeded panel of the rule game (a hidden rule over three integers, one of ten in the game code; six-item quizzes on triples from 0--9; the exam uses held-out instances of the same ten rules, which are textbook hypothesis-testing rules). Re-running the base with a different generation seed gave single-exam SD 0.067 (mean 0.398) and 0.072 (0.445) at 8 episodes, 0.030 (0.416) and 0.044 (0.412) at 40, and 0.033 (0.410) at 80 (population SDs; sample SDs about 7--10\% larger); eighty did not beat forty, so the floor is run-to-run generation variance. Every decision below uses 40 episodes and tolerance 0.03 on ON $-$ OFF. That statistic is a difference of two exams, so its null SD is about $\sqrt{2}$ times the single-exam SD (0.042--0.062), and the tolerance is 0.5--0.7 SD of the decision statistic: decisions near it are coin flips and the $\pm 0.03$ bins sit inside the null. The commit rule as run (ON $\ge$ OFF $-0.03$ plus a paired canary) recorded PREV but did not compare it, so writes could replace better adapters (marked $\dagger$ below). Sequential looks were taken as lineages finished. Earlier 8-episode smoke and scout decisions (6 and 8) are excluded as underpowered.

\paragraph{Design.} One round is 8 classrooms $\times$ 4 lessons $\times$ 2 short episodes from one child snapshot; per lesson the child acts, the parent names one process mistake, the child restates it as a NOTE with scope and applies it to a fresh instance; the lesson's rows enter the corpus only if the post-lesson quiz $\ge$ the pre-lesson quiz. Each round is one v2.1 write (the classroom script passes 3 epochs; the evidence tables say 2; unreconciled against receipts); a lineage is three rounds chained through the latest committed adapter. The comparison arm is self-taught (the 7B child plays parent to itself).

\onoff{} per write on 40-episode exams; ``rej'' $=$ rejected by the harm or canary rule; $\dagger$ $=$ committed below the adapter it replaced; A/B are two machines.
\begin{center}\small
\begin{tabular}{llrrrr}
\toprule
Arm & Lineage & Round 0 & Round 1 & Round 2 & Mean \\
\midrule
14B & s8000, B & $+0.167$ & $+0.008^{\dagger}$ (0.104 below PREV 0.442) & $-0.312$ rej (ON 0.071, collapse) & $-0.046$ \\
14B & s8001, B & $-0.008$ & $+0.012$ & $-0.012^{\dagger}$ (0.108 below PREV 0.450) & $-0.003$ \\
14B & s8000, A & $+0.008$ & $-0.012^{\dagger}$ & $+0.037^{\dagger}$ & $+0.011$ \\
14B & s8001, A & $+0.121$ & $+0.096$ & $+0.017$ & $+0.078$ \\
14B & s8002, B & $-0.017$ & $-0.029$ & $-0.113$ rej & $-0.053$ \\
14B & s8003, B & $+0.100$ & $0.000^{\dagger}$ (0.096 below PREV 0.463) & $-0.100$ rej & $0.000$ \\
14B & s8004, B & $+0.100$ & $+0.054$ & $+0.150$ & $+0.101$ \\
14B & s8005, B & $-0.138$ rej & $-0.150$ rej & $-0.067$ rej & $-0.118$ \\
14B & s8005, A & $-0.104$ rej & $+0.017$ & $-0.238$ rej (ON 0.100, collapse) & $-0.108$ \\
\midrule
self & s8000, B & $-0.025$ & $-0.013$ rej (canary) & $+0.008^{\dagger}$ (ON 0.317 vs PREV 0.375) & $-0.010$ \\
self & s8001, B & $-0.133$ rej & $-0.008$ & $+0.037$ & $-0.035$ \\
self & s8000, A & $-0.075$ rej & $-0.096$ rej & $-0.067$ rej & $-0.079$ \\
self & s8002, B & $-0.096$ rej & $-0.071$ rej & $-0.008$ & $-0.058$ \\
self & s8003, B & $+0.029$ & $+0.071$ & $-0.033$ rej & $+0.022$ \\
self & s8001, A & $-0.075$ rej & $-0.008$ & $-0.067$ rej & $-0.050$ \\
self & s8004, A & $+0.129$ & $+0.021$ & $-0.033$ rej & $+0.039$ \\
self & s8005, B & $-0.083$ rej & $+0.083$ & $+0.017$ & $+0.006$ \\
\bottomrule
\end{tabular}
\end{center}
Tallies: 14B 27 writes, $>+0.03$ in 8, $<-0.03$ in 8, lineage means $\ge 0$ in 4/9; self 24 writes, 4 and 11, means $\ge 0$ in 3/8; seed-to-seed spread (one parented seed $+0.101$; another $-0.118$ and $-0.108$ on two machines) exceeds the arm difference. Harmful writes by round: 14B 2/9, 1/9, 5/9; self 5/8, 2/8, 4/8 (six-lineage interim: 0/6, 0/6, 3/6 and 4/6, 2/6, 3/6). Below-PREV commits: five counted in the notebook (all parented) plus the self s8000 B round-2 write, never tallied; PREV was not logged for every commit, so six is a lower bound. The largest single writes ($+0.167$, $+0.150$, $+0.129$, $+0.121$) are selected maxima from 51 draws whose null maximum is about 0.09--0.13; the $+0.167$ was measured against an OFF exam of 0.342, one of the lowest on record (band mean about 0.41). Single-round cells (seeds 7004--7007; one write from base): 14B $-0.012$, $-0.058$ rej, $+0.058$, $-0.167$ rej; self $-0.138$ rej, $-0.042$ rej, $-0.050$ rej, $-0.021$ rej; 1/8 above $+0.03$, 5/8 below $-0.03$; in-classroom post-minus-pre score positive in 5/8 (no band). Counting the 17 round-0 writes too, single-round writes were above $+0.03$ in 6/25 and below $-0.03$ in 12/25.

\paragraph{Replicate at the tolerance.} One self-taught lineage was started with identical seed and configuration on two machines; only the round-0 writes are a same-configuration pair: $-0.025$ (ON 0.396, OFF 0.421; committed) and $-0.075$ (0.375, 0.450; harm-rejected), each exam within one SD of the other yet flipping the decision. A three-arm parent-size cell on one seed (32B-AWQ $-0.087$ rej, $-0.029$, $-0.054$ rej; 14B $+0.013$, $-0.167$ rej, $+0.075$; self $0.000$, $-0.117$ rej, $-0.100$ rej) had no clean arm and supports no size claim ($n = 1$).

\paragraph{Collapse at s8000 B round 2.} Quiz accuracy 0.071 with the adapter on against 0.383 off (the round-1 adapter it would have replaced scored 0.287). Markers intact (parseable-\texttt{ACT} 0.86 on, 0.91 off; 0/40 episodes began with \texttt{DONE}), but the child acted and wrote less: 5.2 vs.\ 6.3 chunks per episode, 431 vs.\ 728 \texttt{TRY} guesses, 694 vs.\ 870 characters per chunk. One cell, not a mechanism. Its corpus filled 400/400 (256 new experience, 118 remainder, 22 prediction revision, 4 anchors); the child played round 2 under a round-1 adapter itself committed 0.104 below PREV; the round-1 playbook rewrite had been skipped after a context overflow. Leak-fallback replacements of 96 utterances per lineage, where recorded: s8000 B 12; s8001 B 6; s8000 A 4; s8001 A 10; s8002 B 7.

\section{Corpus compositions and schooling renderings (unmeasured)}
\label{app:corpora}

\begin{center}\small
\begin{tabular}{llrrrrrrl}
\toprule
Corpus & Source & Rows & metaflow & opening & contrast & review & full & Status \\
\midrule
v3 & rule-game classroom ledgers & 470 & 206 & 94 & 70 & 65 & 35 & no ritual flags (recipe share 0.26 over the 35 full rows; Jaccard 0.11); development-only, unverified provenance, not leak-scanned \\
v2 & gym-exposed life (structure-heavy) & 370 & 152 & 80 & 61 & 52 & 25 & no flags; quarantined dev-only \\
v1 & gym-exposed life (lived-heavy) & 552 & --- & 200 & 199 & --- & 153 & no flags; quarantined dev-only \\
\bottomrule
\end{tabular}
\end{center}
Renderings are authored by the 32B-AWQ model from the child's stored context (first person, the child's dialect): \emph{metaflow} (pre-action thinking flow, 10--16 lines, ending in one NOTE with plan and stopping rule), \emph{review} (end-of-episode NOTE), \emph{opening} (start of the next episode: RECALL, NOTE, PREDICT), \emph{contrast} (NOTE comparing two programs), \emph{full} (the episode rewritten as PLAN / EXECUTE / REVIEW with every ACT and outcome kept). Instructed rule: never add, remove or change an action or result. Enforced rule (weaker): a non-full rendering is dropped if it contains any ACT line; a full rendering if its ACT lines are not a subset of the original's (47 dropped in v1, 7--10 in v2/v3); removed actions and changed outcomes were not checked. Mixture per episode (v2/v3): metaflow 0.4, review 0.15, opening 0.2, contrast 0.15, full 0.1. Adapters: v2.1 trainer, rank 8, 2 epochs, one run each; v3 final loss 0.675; none run in a life. Per the protocol reviewer's stop order, none may be parented, promoted or reported as a clean child until a fail-closed ancestry check exists. Other corpora: classroom rounds target 400 admitted rows (fills 306/400 and 400/400 logged); recipe screen 130 fixed rows; rank calibration 300, 1{,}200 and 12{,}031 rows; long lives cumulative, deduplicated on the first 160 characters.

\section{Infrastructure, versions and what we tried first (no identifiers)}
\label{app:infra}

Child served by vLLM (16k-token window) with the LoRA hot-swapped per sleep; parents on their own GPU. Two independent machines with 48\,GB-class GPUs ran lives in parallel; a 32B parent in bf16 with two-way tensor parallelism hung three times and was replaced by a 4-bit AWQ build on one GPU. Versions (both machines): child Qwen2.5-7B-Instruct, snapshot \texttt{a09a35458c702b33eeacc393d103063234e8bc28}; torch 2.13.0+cu130; vLLM 0.27.1; PEFT 0.20.0; transformers 5.5.3; GPU driver 580.173.02; CompilerGym version to be recorded (\nm). Each life phase writes a marker file and is idempotent; a gated sleep costs about six minutes more than an ungated one (one timing). Parent utterances are logged with prompt version and model.

\paragraph{What we tried first.} Under an older writer without canonicalization, three seeded lives ran before the recipe was fixed. One was paired-positive at 7 of 8 checkpoints through episode 512 ($+0.006$ to $+0.025$) with no later tally in our evidence sources; a second at 5 of 5 before it was stopped at 344. The third's first sleep captured decorated marker lines and locked that form in; its paired probe had fallen to about $-0.48$ by episode 512 and over its shortened life (stopped near 704) it averaged $-0.219$ with 7 harmful pairs of 10. That bifurcation produced the canonicalization step and the real-context canary. An earlier pilot also leaked the probe set through the identifier prefix (\S\ref{sec:measure}). Those failures are the recipe reported here.

\section{Additional related work: sleep consolidation}
\label{app:related}

Complementary learning systems \citep{mcclelland1995cls} justify a fast episodic store (the ledger) and slow interleaved consolidation (the write); active system consolidation \citep{diekelmann2010sleep} licenses repeated small cycles; synaptic homeostasis \citep{tononi2014sleep} predicts survival of the consistent, the analogue of our amplifier reading; replay is biased \citep{buzsaki2015swr}; consolidation must be regulated \citep{sun2023organizing}; targeted reactivation \citep{oudiette2013tmr} shows selection is steerable from outside, the nearest analogue of a parent who never writes the corpus. Wake--sleep \citep{hinton1995wakesleep} and generative replay \citep{shin2017replay,arani2022clser} are the machine-learning ancestors; we keep a verbatim ledger because generated replay drifts. Self-modification with dreaming curricula \citep{behrouz2026sleep} is a training paradigm over static corpora; ``dream'' in the world-model sense \citep{hafner2023dreamerv3,nottingham2023deckard} is imagined rollouts, not context management. \citet{zahavy2026jump} argues that models lack abduction, the invention of an explanation for a surprise; the loop here trains its ingredients where checkable evidence exists.

\end{document}
